% Standalone section. Requires amsmath, amssymb, amsthm and hyperref.
% The parent article should provide theorem, proposition, corollary, remark.
% Bibliography keys: dhr20261010, dlmfgauss, dlmfgammaratio, dlmfhurwitz, dlmfpolylog.
\section{The next two-particle resonance: overlap subtraction and exact constants}
\label{sec:second-resonance}

The determinantal harmonic-resonance report explicitly proposed the next
two-particle resonance as the first test of a higher-pole subtraction
\cite[Section~9]{incoming-determinantal}. We resolve that target for arbitrary positive
shape parameters, and by joint holomorphy for their right half-planes.
The result consists of an explicit normally convergent subtraction, its
complete spectral derivative calculus, and a finite-cutoff constant proved
independently from an exact discrete Pearson identity. These statements
are proved here; no claim of exhaustive literature novelty is made.

\subsection{Definitions and a finite identity}
Let $\Re a,\Re c>0$ and put
\begin{align}
 w_n(u)&=\frac{\Gamma(n+a)\Gamma(n+c)}
 {\Gamma(n+a+c+1+u)\Gamma(n+1)},\label{sr:eq:weight}\\
 S_N(u)&=\sum_{0\le m<n\le N}(n-m)^2w_m(u)w_n(u),\qquad
 K_n(u)=\sum_{0\le m<n}(n-m)^2w_m(u)w_n(u).
 \label{sr:eq:cutoff}
\end{align}
Thus $S_N=\sum_{n=0}^N K_n$ and $K_0=0$. There is no factor $2!$.
Use the abbreviations
\begin{gather}
 \mu=\frac1{ac},\qquad h=\frac1a+\frac1c,\qquad q=1+h,
 \qquad s=ac+a+c,\label{sr:eq:abbreviations}\\
 \rho=\psi(1)-\psi(a+1)-\psi(c+1),\qquad
 \tau=\psi'(1)-\psi'(a+1)-\psi'(c+1),\qquad r_2=\rho^2+\tau.
 \label{sr:eq:rho}
\end{gather}
The letter $q$ in this section is a scalar abbreviation, not a
cyclotomic level or an integration order.

\begin{proposition}[Exact Pearson reduction]\label{sr:prop:pearson}
Write $M_j(N,u)=\sum_{n=0}^N n^jw_n(u)$ and
$t_n(u)=n(n+a+c+u)w_n(u)$. For every $N\ge0$,
\begin{align}
 acM_0(N,u)-uM_1(N,u)&=t_{N+1}(u),\label{sr:eq:pearson0}\\
 (1-u)M_2(N,u)+(s+u)M_1(N,u)&=Nt_{N+1}(u).
 \label{sr:eq:pearson1}
\end{align}
Consequently, at $u=0$,
\begin{align}
 acM_0(N,0)&=(N+1)(N+1+a+c)w_{N+1}(0),\label{sr:eq:M0exact}\\
 S_N(0)&=acN M_0(N,0)^2-sM_0(N,0)M_1(N,0)-M_1(N,0)^2.
 \label{sr:eq:finite-det}
\end{align}
\end{proposition}
\begin{proof}
The exact weight ratio is
\[
 \frac{w_{n+1}(u)}{w_n(u)}
 =\frac{(n+a)(n+c)}{(n+a+c+1+u)(n+1)}.
\]
It gives
\[
 t_{n+1}-t_n=(ac-un)w_n,
 \qquad
 (n+1)t_{n+1}-nt_n=((1-u)n^2+sn+ac)w_n.
\]
Sum these identities, use $t_0=0$, and eliminate $acM_0$ to obtain
\eqref{sr:eq:pearson0}--\eqref{sr:eq:pearson1}. The elementary identity
$S_N=M_0M_2-M_1^2$ follows by expanding $(n-m)^2$ and pairing the two
strict orderings. Substitution at zero proves \eqref{sr:eq:finite-det}.
\end{proof}

\subsection{The Gamma product and its two different singular scales}
Define holomorphic germs
\begin{align}
 R(u)&=\frac{\Gamma(a+1)\Gamma(c+1)\Gamma(1+u)}
 {\Gamma(a+1+u)\Gamma(c+1+u)},& A(u)&=\mu R(u),\label{sr:eq:R}\\
 \beta(u)&=-s-(a+c+\tfrac12)u-\tfrac12u^2,
 &V(u)&=\frac1{1-u^2},\label{sr:eq:beta}\\
 P(u)&=(1+hu+\mu u^2)R(u)^2,
 &T(u;\alpha)&=A(u)(\beta(u)+\alpha u).
 \label{sr:eq:PT}
\end{align}
In particular $R(0)=1$, $R'(0)=\rho$, and $R''(0)=r_2$.
For $\Re u>1$, the partition converges absolutely, and Gauss summation
\cite[15.4.20]{dlmfgauss} gives
\begin{equation}\label{sr:eq:partition-product}
 \mathcal Z(u):=\sum_{n\ge0}K_n(u)
 =\frac{P(u)}{u^2(u-1)}.
\end{equation}
For completeness, its three moment sums in this domain are
\[
 M_0(\infty,u)=\mu R(u),\qquad
 M_1(\infty,u)=\frac{R(u)}u,\qquad
 M_2(\infty,u)=\frac{R(u)(s+u)}{u(u-1)}.
\]
Their determinant is \eqref{sr:eq:partition-product}. Thus the product
needs only ordinary Gauss summation, not an assumed higher-dimensional
Selberg formula.

Fix $\Re\alpha>0$ and write $y=n+\alpha$, with the logarithm on the
right half-plane. Define the \emph{single grouped counterterm}
\begin{align}
 Q_n(u;\alpha)={}&A(u)y^{-u}+T(u;\alpha)y^{-1-u}\notag\\
 &+\frac2u\left\{V(u)y^{-1-2u}-R(u)y^{-1-u}\right\}.
 \label{sr:eq:Q}
\end{align}
The last line is holomorphic at $u=0$ because its numerator vanishes
there. The exponent $1+2u$ is essential: it accounts for the region in
which the inner summation index also becomes large. Replacing it by
$1+u$ loses the logarithmic overlap term.

\begin{theorem}[Normally convergent second-resonance germ]
\label{sr:thm:germ}
For $\Re a,\Re c,\Re\alpha>0$, the series
\begin{equation}\label{sr:eq:D}
 \mathcal D(u;\alpha)=\sum_{n\ge0}\{K_n(u)-Q_n(u;\alpha)\}
\end{equation}
converges absolutely and locally normally for $u$ in a neighborhood of
zero. It is jointly holomorphic in these parameters, and every fixed
mixed parameter derivative can be taken termwise. Its exact value is
\begin{align}
 \mathcal D(u;\alpha)={}&\frac{P(u)}{u^2(u-1)}-A(u)\zeta(u,\alpha)
 -T(u;\alpha)\zeta(1+u,\alpha)\notag\\
 &+\frac{2R(u)}u\zeta(1+u,\alpha)
 -\frac{2V(u)}u\zeta(1+2u,\alpha).
 \label{sr:eq:master}
\end{align}
All apparent poles at $u=0$ cancel in this complete expression.
\end{theorem}
\begin{proof}
Work on compact shape-parameter sets and choose $0<\delta<1/4$.
The locally uniform Gamma-ratio expansion \cite[5.11.13]{dlmfgammaratio}
is
\begin{equation}\label{sr:eq:w-asymptotic}
 w_n(u)=n^{-2-u}\left(1+\frac{\beta(u)}n+O(n^{-2})\right),
 \qquad |u|\le\delta.
\end{equation}
Gauss summation gives $\sum_{m\ge0}w_m(u)=A(u)$ already in
$\Re u>-1$. Hence, uniformly on the chosen compact sets,
\begin{equation}\label{sr:eq:lower-M0}
 \sum_{m<n}w_m(u)=A(u)-\frac{n^{-1-u}}{1+u}+O(n^{-2+\delta}).
\end{equation}
For the first moment, the series
$\sum_{m\ge1}\{mw_m(u)-m^{-1-u}\}$ converges locally normally for
$\Re u>-1$. On $\Re u>0$ its value is $R(u)/u-\zeta(1+u)$, so that
identity continues through zero. Its tail, together with the leading
Hurwitz-zeta tail, gives the \emph{holomorphic} estimate
\begin{equation}\label{sr:eq:lower-M1}
 \sum_{m<n}mw_m(u)=\frac{R(u)-n^{-u}}u+O(n^{-1+\delta}).
\end{equation}
The quotient has its removable value at zero. Finally, summing
$m^2w_m(u)=m^{-u}+O(m^{-1+\delta})$ gives
\begin{equation}\label{sr:eq:lower-M2}
 \sum_{m<n}m^2w_m(u)=\frac{n^{1-u}}{1-u}
 +O(n^\delta\log(n+1)).
\end{equation}
Insert these estimates in
\[
 K_n=w_n\left(n^2\sum_{m<n}w_m-2n\sum_{m<n}mw_m
                         +\sum_{m<n}m^2w_m\right).
\]
The coefficient of $n^{-1-2u}$ is
\[
 -\frac1{1+u}+\frac2u+\frac1{1-u}
 =\frac{2}{u(1-u^2)}.
\]
The resulting paired expression is holomorphic at zero, and
\[
 K_n=A(u)n^{-u}+A(u)\beta(u)n^{-1-u}
 +\frac2u\{V(u)n^{-1-2u}-R(u)n^{-1-u}\}
 +O(n^{-2+2\delta}\log(n+1)).
\]
Changing $n$ to $y=n+\alpha$ in the leading scale adds
$\alpha uA(u)y^{-1-u}$; the remaining changes are summable uniformly.
This proves \eqref{sr:eq:D}, since the displayed error is summable.
Cauchy's integral formula on slightly larger compact parameter sets
proves the assertions for every fixed derivative.

To identify the sum, the same estimates hold on compact subsets of
$-\delta<\Re u<1+\delta$, excluding $u=1$. In particular,
\eqref{sr:eq:lower-M2} still holds there: summing $m^{-u}$ gives
$n^{1-u}/(1-u)$ plus a locally bounded constant and
$O(n^\delta)$, while the partial-sum estimate for
$m^2w_m-m^{-u}$ is $O(n^\delta\log(n+1))$. Constants are uniform
on compact sets avoiding $u=1$. Thus the subtracted remainder remains
summable throughout this strip. Both complete expressions are
holomorphic on its connected punctured domain, including their
removable values at zero.
On $1<\Re u<1+\delta$ all original sums and counterterms converge
separately, and their difference is \eqref{sr:eq:master} by
\eqref{sr:eq:partition-product}. The identity theorem gives the
formula near zero. Its removability there was established by the
normally convergent left side, independently of formal Laurent algebra.
\end{proof}

\subsection{Exact hard-cutoff constant and spectral comparison}
Set
\begin{equation}\label{sr:eq:lambda-C}
 \lambda=-q-2\rho,\qquad C=-\rho^2-q\rho-2.
\end{equation}
\begin{theorem}[Closed constant at the next resonance]
\label{sr:thm:constant}
As $N\to\infty$,
\begin{equation}\label{sr:eq:cutoff-asymptotic}
 \boxed{\displaystyle
 S_N(0)=\mu N-\log^2N+\lambda\log N+C+o(1).}
\end{equation}
The error is $O(\log N/N)$ for fixed shape parameters, locally uniformly
on their right half-planes. The spectral Laurent expansion is instead
\begin{equation}\label{sr:eq:spectral-laurent}
 \mathcal Z(u)=-\frac1{u^2}+\frac\lambda u+F+O(u),\qquad
 F=-\mu-q-2q\rho-2\rho^2-\tau.
\end{equation}
Thus the exact correction between the two constants is
\begin{equation}\label{sr:eq:finitepart-correction}
 F-C=-\mu-q+2-q\rho-\rho^2-\tau.
\end{equation}
\end{theorem}
\begin{proof}
At zero, \eqref{sr:eq:lower-M0}--\eqref{sr:eq:lower-M1} give
\[
 M_0(N,0)=\mu-\frac1N+O(N^{-2}),\qquad
 M_1(N,0)=\log N+\rho+O(N^{-1}).
\]
Substituting them in the exact finite identity
\eqref{sr:eq:finite-det} gives
\[
 acNM_0^2=\mu N-2+O(N^{-1}),\quad
 sM_0M_1=q(\log N+\rho)+O(\log N/N),
\]
and the asserted cutoff formula follows. This proof does not identify
a spectral finite part with a cutoff constant. Separately expanding
$R(u)=1+\rho u+(\rho^2+\tau)u^2/2+O(u^3)$ in
\eqref{sr:eq:partition-product} proves \eqref{sr:eq:spectral-laurent}
and \eqref{sr:eq:finitepart-correction}.
\end{proof}

With the Stieltjes convention
\[
 \zeta(1+u,\alpha)=\frac1u+
 \sum_{k\ge0}\frac{(-1)^k}{k!}\gamma_k(\alpha)u^k,
 \qquad \gamma_0(\alpha)=-\psi(\alpha),
\]
the constant term of Theorem~\ref{sr:thm:germ} is the ordinary
single-bracket identity
\begin{equation}\label{sr:eq:D0}
 \boxed{\displaystyle
 \sum_{n\ge0}\left\{K_n(0)-\mu-
 \frac{\lambda-2\log(n+\alpha)}{n+\alpha}\right\}
 =C-\mu-\lambda\gamma_0(\alpha)+2\gamma_1(\alpha).}
\end{equation}
Its agreement with \eqref{sr:eq:cutoff-asymptotic} also follows by
summing the elementary counterterms before taking the limit. The
coefficient $2\gamma_1$ comes from the doubled exponent $2u$.

\subsection{Every derivative and an explicit first row}
Write $R_j=R^{(j)}(0)$, $T_j=\partial_u^jT(0;\alpha)$ and
$V_j=V^{(j)}(0)$, so $V_j=j!$ for even $j$ and $V_j=0$ for odd $j$.
Let
\begin{equation}\label{sr:eq:N-germ}
 \mathcal N(u)=\frac{P(u)}{u-1}+2R(u)-V(u)-uT(u;\alpha).
\end{equation}
Direct substitution gives $\mathcal N(0)=\mathcal N'(0)=0$.
Thus $\mathcal N(u)/u^2$ is an explicitly given holomorphic
Gamma/polygamma germ, not an additional unspecified constant.

\begin{corollary}[Finite formula for every spectral row]
\label{sr:cor:alljets}
For every $k\ge0$, put $D_k(\alpha)=\partial_u^k\mathcal D(0;\alpha)$.
Then
\begin{align}
 D_k(\alpha)={}&
 \frac{\mathcal N^{(k+2)}(0)}{(k+1)(k+2)}
 -\mu\sum_{j=0}^k\binom kj R_j\zeta^{(k-j)}(0,\alpha)\notag\\
 &-\sum_{j=0}^k\binom kj T_j(-1)^{k-j}\gamma_{k-j}(\alpha)\notag\\
 &+\frac2{k+1}\sum_{j=0}^{k+1}\binom{k+1}j
 \bigl(R_j-2^{k+1-j}V_j\bigr)(-1)^{k+1-j}
 \gamma_{k+1-j}(\alpha).
 \label{sr:eq:alljets}
\end{align}
Here $\zeta^{(j)}$ denotes differentiation in the order variable.
The normally convergent left side is
$\sum_{n\ge0}\{K_n^{(k)}(0)-Q_n^{(k)}(0;\alpha)\}$. If
$L_n=\log(n+\alpha)$, every counterterm is the finite polynomial
\begin{align}
 Q_n^{(k)}(0;\alpha)={}&\mu\sum_{j=0}^k\binom kjR_j(-L_n)^{k-j}
 +\frac1{n+\alpha}\sum_{j=0}^k\binom kjT_j(-L_n)^{k-j}\notag\\
 &+\frac2{(k+1)(n+\alpha)}
 \sum_{j=0}^{k+1}\binom{k+1}j
 \left\{V_j(-2L_n)^{k+1-j}-R_j(-L_n)^{k+1-j}\right\}.
 \label{sr:eq:counterterm-jets}
\end{align}
\end{corollary}
\begin{proof}
Put $J(u;\alpha)=\zeta(1+u,\alpha)-1/u$. Rearranging
\eqref{sr:eq:master} gives the manifestly regular expression
\[
 \mathcal D=\frac{\mathcal N(u)}{u^2}-\mu R(u)\zeta(u,\alpha)
 -T(u;\alpha)J(u;\alpha)
 +\frac2u\{R(u)J(u;\alpha)-V(u)J(2u;\alpha)\}.
\]
Apply the product rule. Division of a germ vanishing to order $m$
converts its $(k+m)$th derivative to its $k$th derivative divided by
$(k+1)\cdots(k+m)$; this gives the two denominators in
\eqref{sr:eq:alljets}. The same argument applied to \eqref{sr:eq:Q}
gives \eqref{sr:eq:counterterm-jets}. Normal convergence was already
proved in Theorem~\ref{sr:thm:germ}.
\end{proof}

All coefficients in this formula are finite polygamma polynomials.
For example, the complete exponential Bell polynomials give
\[
 R_j=\operatorname{Bell}_j(\ell_1,\ldots,\ell_j),\quad
 \ell_j=\psi^{(j-1)}(1)-\psi^{(j-1)}(a+1)-\psi^{(j-1)}(c+1).
\]
The other coefficients follow by multiplication by the explicit
polynomials and rational functions in \eqref{sr:eq:beta}--\eqref{sr:eq:PT}.
In fact the apparently required derivative $R_{k+2}$ always cancels.
To see this without coefficient cancellation, put
$W(u)=(R(u)-1)/u$ and $d=\mu(\alpha-a-c-\tfrac12)$. Then
\begin{equation}\label{sr:eq:lower-order-G}
 \frac{\mathcal N(u)}{u^2}=-W(u)^2-V(u)
 +R(u)\left\{\frac{(q+\mu u)W(u)+q+\mu}{u-1}
                    -d+\frac\mu2u\right\}.
\end{equation}
Its $k$th derivative uses only $R_0,\ldots,R_{k+1}$.
Thus the $k$th identity needs shape polygammas only through order $k$,
while its Stieltjes constants run through order $k+1$.
The summand itself has the equally explicit expression
\begin{equation}\label{sr:eq:K-bell}
 K_n^{(k)}(0)=\sum_{m<n}(n-m)^2w_m(0)w_n(0)
 \operatorname{Bell}_k(x_1(m,n),\ldots,x_k(m,n)),
\end{equation}
where
\[
 x_j(m,n)=-\psi^{(j-1)}(m+a+c+1)-\psi^{(j-1)}(n+a+c+1).
\]
At positive integer parameters these are ordinary generalized harmonic
numbers and zeta constants; at arbitrary shifts they are their
polygamma continuations.

The first derivative can be printed without Bell notation. Define
\begin{align}
 B(\alpha)&=r_2+q\rho+q-3+\mu(\tfrac12-\alpha),\label{sr:eq:B}\\
 C_1&=-\rho r_2-q\rho^2-\frac q2r_2
 -(2\mu+q+1)\rho-q-\frac\mu2.
 \label{sr:eq:C1}
\end{align}
Then
\begin{align}
 Q_n'(0;\alpha)&=\mu(\rho-L_n)
 +\frac{3L_n^2-\lambda L_n-B(\alpha)}{n+\alpha},
 \label{sr:eq:Q1}\\
 D_1(\alpha)&=C_1-\mu\zeta'(0,\alpha)
 +B(\alpha)\gamma_0(\alpha)+\lambda\gamma_1(\alpha)
 -3\gamma_2(\alpha).
 \label{sr:eq:D1}
\end{align}
Here $K_n'(0)$ is minus the defining double summand multiplied by
$\psi(m+a+c+1)+\psi(n+a+c+1)$. The term
$\zeta'(0,\alpha)=\log\Gamma(\alpha)-\tfrac12\log(2\pi)$
is classical \cite[25.11.18]{dlmf-hurwitz}. Thus this row is an explicit
ordinary convergent identity connecting a polygamma-weighted double
sum, log-Gamma, and the first three generalized Stieltjes constants.

For comparison, the corresponding unsubtracted finite cutoff has
the full first-derivative expansion
\begin{align}
 S_N'(0)={}&-\mu N\log N+\mu(\rho+1)N+\log^3N
 -\frac\lambda2\log^2N\notag\\
 &-(r_2+q\rho+q-3+\mu)\log N+C_1+\mu\rho+o(1).
 \label{sr:eq:cutoff-derivative}
\end{align}
Indeed take $\alpha=1$ in \eqref{sr:eq:Q1}, sum through $N$,
use $\sum_{j=1}^{N+1}\log j=\log\Gamma(N+2)$ and the defining
Stieltjes limits, and apply Stirling's formula. The constants
$\gamma_0,\gamma_1,\gamma_2$ cancel against \eqref{sr:eq:D1}; the
$\log(2\pi)$ terms cancel as well.

\subsection{A uniform bridge from spectral germs to finite cutoffs}
For any fixed finite derivative order, one can obtain the full cutoff
expansion without making new interchanges. If $Y=N+1+\alpha$, the
normally convergent tail in Theorem~\ref{sr:thm:germ} gives
\begin{align}
 S_N(u)={}&\mathcal Z(u)-A(u)\zeta(u,Y)
 -\left\{T(u;\alpha)-\frac{2R(u)}u\right\}\zeta(1+u,Y)\notag\\
 &-\frac{2V(u)}u\zeta(1+2u,Y)
 +O(N^{-1+2\delta}\log(N+1)).
 \label{sr:eq:uniform-cutoff}
\end{align}
The complete expression has removable poles at zero. The error is
locally uniform for $|u|\le\delta<1/4$, after shrinking a surrounding
disk if needed for derivatives. Thus every fixed Taylor coefficient
may be extracted, with a remainder tending to zero. The standard
large-$Y$ Hurwitz expansion then generates every $N\log^jN$ term,
every logarithmic counterterm, and its constant. The overlap exponent
$2u$ controls the higher logarithmic powers.

\subsection{Shift differentiation and an anchored primitive}
The subtraction point is auxiliary; it does not move the original
particle grid. Nevertheless \eqref{sr:eq:D0} has the exact laws
\begin{align}
 D_0(\alpha+1)-D_0(\alpha)
 &=\frac{\lambda-2\log\alpha}{\alpha},\label{sr:eq:alpha-shift}\\
 \frac{d}{d\alpha}D_0(\alpha)
 &=(\lambda+2)\zeta(2,\alpha)+2\zeta'(2,\alpha).
 \label{sr:eq:alpha-derivative}
\end{align}
These follow from the Hurwitz shift and derivative identities
\cite[25.11.3 and 25.11.17]{dlmf-hurwitz}. For any anchor with
$\Re\alpha_0>0$, a fixed-constant primitive is
\begin{align}
 \int_{\alpha_0}^{\alpha}D_0(t)\,dt
 ={}&(C-\mu)(\alpha-\alpha_0)
 +\lambda\log\frac{\Gamma(\alpha)}{\Gamma(\alpha_0)}\notag\\
 &+\zeta''(0,\alpha)-\zeta''(0,\alpha_0).
 \label{sr:eq:primitive}
\end{align}
The logarithm is the analytic log-Gamma on the right half-plane.
To verify the last term, differentiate
$\partial_\alpha\zeta(u,\alpha)=-u\zeta(1+u,\alpha)$ twice in
$u$ at zero; the result is
$\partial_\alpha\zeta''(0,\alpha)=2\gamma_1(\alpha)$.
This also explains why an anchored primitive involves the second
order derivative of Hurwitz zeta rather than an unspecified new constant.

\subsection{Two normalization checks and a rational specialization}
If every weight is multiplied by an arbitrary holomorphic $L(u)$
with $L(0)=1$, then $K_n$, $Q_n$, $\mathcal D$, and $\mathcal Z$
are all multiplied by $L(u)^2$. Therefore the cutoff at $u=0$ and
its constant $C$ are unchanged. Writing $l_1=L'(0)$ and
$l_2=L''(0)$, the spectral finite part becomes
\begin{equation}\label{sr:eq:normalization}
 \widetilde F=F+2l_1\lambda-l_1^2-l_2,
\end{equation}
whereas $\widetilde D_1=D_1+2l_1D_0$. These identities explicitly
track both amplitude derivatives relevant to a double pole.

At $a=c=1$ the weight is $w_n(0)=1/((n+1)(n+2))$ and
\[
 \mu=1,\quad \rho=\gamma-2,\quad \lambda=1-2\gamma,
 \quad C=\gamma-\gamma^2.
\]
An exact largest-index formula, obtained from the finite moments, is
\begin{equation}\label{sr:eq:rational-K}
 K_{v-1}(0)=1+\frac1{v+1}-\frac{H_v}{v}-\frac{H_v}{v+1},
 \qquad v\ge1.
\end{equation}
Its first value is zero, as required by the empty inner sum. Thus
\begin{equation}\label{sr:eq:integer-cutoff}
 \sum_{0\le m<n\le N}\frac{(n-m)^2}{(m+1)(m+2)(n+1)(n+2)}
 =N-\log^2N+(1-2\gamma)\log N+\gamma-\gamma^2+o(1).
\end{equation}
For the unnormalized weights the spectral finite part is
$F=\zeta(2)-2+2\gamma-2\gamma^2$. With the different normalization
$L(u)=\Gamma(1+u)$, the entire spectral product simplifies to
\[
 L(u)^2\mathcal Z(u)=\frac1{u^2(u-1)(1+u)^2},
\]
whose finite part is $-2$. Both have exactly the same finite-cutoff
constant $\gamma-\gamma^2$. This is a concrete second-order check
of \eqref{sr:eq:normalization}.

The specialization of \eqref{sr:eq:D0} at $a=c=\alpha=1$ yields
\[
 \sum_{v\ge1}\left\{
 \frac1{v+1}-\frac{H_v}{v}-\frac{H_v}{v+1}
 -\frac{1-2\gamma-2\log v}{v}\right\}
 =\gamma^2-1+2\gamma_1.
\]
This scalar specialization reduces to a classical harmonic--Stieltjes
relation; its novelty is not asserted. The general-shape theorem, the
uniform overlap subtraction, and the all-order cutoff comparison are
the substantive continuation of the proposed resonance problem.

\subsection{A functional polylogarithm identity and convergent endpoints}
The rational specialization has an exact disk identity as well as a
cutoff expansion. Define, for $|z|<1$,
\[
 \mathcal G(z)=\sum_{0\le m<n}
 \frac{(n-m)^2z^{n+1}}{(m+1)(m+2)(n+1)(n+2)}.
\]
The finite formula \eqref{sr:eq:rational-K} gives
\begin{equation}\label{sr:eq:G-polylog}
 \boxed{\displaystyle
 \mathcal G(z)=\frac z{1-z}-1-\frac{\log(1-z)}z
 -\operatorname{Li}_2(z)-\frac12\left(1+\frac1z\right)\log^2(1-z).}
\end{equation}
All apparent singularities at zero are removable. The first term is
$z^2/12$, corresponding to the configuration $(m,n)=(0,1)$.
To prove the formula, sum the four terms of
\eqref{sr:eq:rational-K}; the two harmonic transforms follow by
integrating $\sum_{v\ge1}H_vz^v=-\log(1-z)/(1-z)$ once, either with
$dz$ or $dz/z$. Thus this is a functional certificate before any
special-value substitution.

At $z=1/2$, the classical dilogarithm reflection identity
\cite[25.12.6]{dlmfpolylog} gives
\begin{equation}\label{sr:eq:half-polylog}
 \sum_{0\le m<n}\frac{(n-m)^2 2^{-(n+1)}}
 {(m+1)(m+2)(n+1)(n+2)}
 =2\log2-\log^22-\frac{\pi^2}{12}.
\end{equation}
For $p\ge2$, the ordinary convergent endpoint
\[
 J_p=\sum_{0\le m<n}\frac{(n-m)^2}
 {(m+1)(m+2)(n+1)(n+2)(n+1)^p}
\]
has a finite all-$p$ reduction. To distinguish multiple zeta values
from Hurwitz zeta, write
$\zeta_{\mathrm{MZ}}(j,1)=\sum_{v>m\ge1}v^{-j}m^{-1}$.
Set $E_j=\zeta_{\mathrm{MZ}}(j,1)+\zeta(j+1)$ for
$j\ge2$; Euler's height-one formula reduces every $E_j$ to products
of at most two ordinary zeta values \cite[Section~6]{incoming-determinantal}.
Then
\begin{equation}\label{sr:eq:allp-endpoint}
 J_p=\zeta(p)-E_{p+1}
 +\sum_{j=2}^{p}(-1)^{p-j}\bigl(\zeta(j)-E_j\bigr)
 +(-1)^{p-1}(1-\zeta(2)).
\end{equation}
Indeed partial fractions give
\begin{align*}
 \sum_{v\ge1}\frac1{v^p(v+1)}
 &=\sum_{j=2}^p(-1)^{p-j}\zeta(j)+(-1)^{p-1},\\
 \sum_{v\ge1}\frac{H_v}{v^p(v+1)}
 &=\sum_{j=2}^p(-1)^{p-j}E_j+(-1)^{p-1}\zeta(2);
\end{align*}
substitute them in \eqref{sr:eq:rational-K}. In particular,
\begin{align}
 J_2&=3\zeta(2)-2\zeta(3)-\tfrac54\zeta(4)-1,\label{sr:eq:J2}\\
 J_3&=1-2\zeta(2)+4\zeta(3)-\tfrac54\zeta(4)
       +\zeta(2)\zeta(3)-3\zeta(5).
 \label{sr:eq:J3}
\end{align}
The primitive of order one has a simple explicitly subtracted endpoint:
\begin{equation}\label{sr:eq:subtracted-J1}
 \sum_{v\ge1}\frac{K_{v-1}(0)-1}{v}
 =1-\zeta(2)-2\zeta(3).
\end{equation}
Its convergence follows from $K_{v-1}(0)-1=O(\log v/v)$.

\subsection{What this resolves, and what remains}
The next $r=2$ resonance requested in the earlier report is resolved:
two moment divergences require three asymptotic counterterm scales,
the overlap term is explicit, normal convergence is proved, and the
cutoff/spectral constants are compared including normalization changes.
The result does not establish a general higher-resonance theorem at
arbitrary particle number. Natural next problems are to identify the
escape strata at $r\ge3$, formulate an inclusion--exclusion subtraction
for simultaneous escape, and discover finite Pearson or contiguous
relations that reduce the number of independent divergent moments.
Another precise question is whether the order saving proved in
\eqref{sr:eq:lower-order-G} survives at higher particle number, where
several nontrivial residue amplitudes replace the single $R(u)$.
